\section{Related Work}
\label{sec:related-work}

Stale coherence sits in a gap between two literatures: long-context and
memory benchmarks (\S\ref{sec:rw-memory}), which measure verbal or action
recall but not their joint behaviour on tick-continuous context, and
production-agent failure observations (\S\ref{sec:rw-production}), which
anecdotally report the phenomenon but do not measure it. We organize this
section accordingly.

\subsection{Long-Context and Memory Benchmarks}
\label{sec:rw-memory}

\paragraph{Verbal-recall axis.}
The dominant long-context memory regime probes verbal retrieval:
\textbf{LongMemEval}~\citep{Wu2025LongMemEval} curates multi-session
conversational memory queries; \textbf{RULER}~\citep{Hsieh2024RULER}
synthesizes needle-in-a-haystack and aggregation tasks across context lengths;
\textbf{Lost-in-the-Middle}~\citep{Liu2024LostInMiddle} characterizes the
U-shaped position bias of retrieval over long inputs;
\textbf{InfiniteBench}~\citep{Zhang2024InfiniteBench} pushes the same probe
beyond $10^5$-token contexts. HELM's recent long-context
battery~\citep{Stanford2025HELMLong} aggregates these into an omnibus
verbal-recall score. In each case the figure of merit is ``did the model
retrieve the fact?''; none ask ``did it \emph{act} on it?''. As a result a
model that hallucinates an anchor in its summary while ignoring it in its
directives is indistinguishable from one that genuinely sustains commitment.

\paragraph{Action-grounding axis (session-discrete).}
A second strand evaluates whether the agent \emph{acts} on remembered facts:
\textbf{MemoryArena}~\citep{He2026MemoryArena} probes action grounding by
issuing explicit downstream queries that require recalled facts;
\textbf{MemoryAgentBench}~\citep{Hu2025MemoryAgentBench} adds selective
forgetting and structured tasks. These are the only benchmarks that touch the
action axis at all, but both are \emph{session-discrete}: each task has
explicit boundaries, the agent receives an explicit query that re-cues the
fact, and the action set is small. Stale coherence, in contrast, requires
\emph{tick-continuous} naturally-accumulated context: the agent emits
structured directives continuously, never receives an explicit query about
the anchor, and decay is measured by comparison across time rather than by
query-response matching. The two regimes probe different failure modes ---
session-discrete action grounding asks whether the agent retrieves on demand;
stale coherence asks whether it sustains commitment when no demand is issued.

\paragraph{Theoretical underpinning.}
\citet{Sinha2025IllusionDiminishingReturns} characterizes long-context recall
as following a diminishing-returns curve: per-token information yield decays
sub-linearly with effective context length. We adopt this geometry as the
expected functional form for both axes of the dual probe and report decay
profiles for $\Delta(t)$ rather than scalar end-of-episode scores. The
contribution of this paper is the operational counterpart of that theory for
the \emph{action} axis, and the demonstration that the gap between the two
axes is itself a measurable independent phenomenon.

\paragraph{Position.}
Stale coherence is the dissociation visible only when both axes are measured
jointly on tick-continuous context. No existing benchmark probes both axes
simultaneously, and no existing benchmark runs long enough on
naturally-accumulated agent context for the dissociation to emerge. The
methodological apparatus of \S\ref{sec:metrics} --- dual probe, anchor
injection at $t = 0$, 8-sim-hour episodes --- is built specifically to close
this gap.

\subsection{Production Agent Failure Observations}
\label{sec:rw-production}

The phenomenon we name has been reported anecdotally in production agent
deployments. \textbf{FAIR Cicero}~\citep{FAIR2022Cicero}, the Diplomacy
negotiation agent, informally documents cases in which the system continued
to declare strategic intent (``defend territory~$X$'') in its dialogue
generations while its unit-movement directives shifted away from that
territory --- a textbook instance of verbal-action dissociation in a
production system. Anthropic's multi-agent research
system~\citep{AnthropicMultiAgent2025} reports inter-channel coherence loss
in long-running sessions of its production multi-agent deployment, including
cases where summary-channel outputs continued to assert commitments that the
acting-channel outputs no longer reflected.

These reports are anecdotal, qualitative, and unreproducible: they describe
the failure mode but do not measure it, isolate it from raw recall, or
quantify its prevalence across model families. The contribution of this
paper, relative to these observations, is to take the same phenomenon and
deliver controlled, quantified, reproducible measurement: a single anchor,
two paired probes, an 8-sim-hour deterministic substrate, a single
pre-registered hypothesis, and three model families.
